
\documentclass[11pt]{article}
\usepackage{amsmath,amssymb,geometry,booktabs}
\geometry{margin=1in}

\title{Five-State Estimation Model for a Differential Drive Robot}
\author{Dash Robot Miniature AV Stack -- Phase 1 (Open Loop Estimation)}
\date{}

\begin{document}
\maketitle

\section*{Overview}

This document defines a five–state kinematic model suitable for a miniature
autonomous vehicle (AV) estimation stack. The robot streams wheel encoder data
to a host computer, which runs an Extended Kalman Filter (EKF) to estimate the
robot state.

The estimator operates in discrete time with timestep $\Delta t$.

\section*{State Vector}

The state vector is

\begin{equation}
x_k =
\begin{bmatrix}
p_x \\
p_y \\
\theta \\
v \\
\omega
\end{bmatrix}
\in \mathbb{R}^5
\end{equation}

\begin{center}
\begin{tabular}{lll}
\toprule
State & Description & Units \\
\midrule
$p_x$ & Position in world $x$ direction & meters \\
$p_y$ & Position in world $y$ direction & meters \\
$\theta$ & Heading (yaw) angle & radians \\
$v$ & Forward velocity & m/s \\
$\omega$ & Yaw rate & rad/s \\
\bottomrule
\end{tabular}
\end{center}

\section*{Process Dynamics}

The robot follows differential–drive kinematics.

\begin{align}
p_{x,k+1} &= p_{x,k} + v_k \cos(\theta_k)\,\Delta t \\
p_{y,k+1} &= p_{y,k} + v_k \sin(\theta_k)\,\Delta t \\
\theta_{k+1} &= \theta_k + \omega_k \Delta t \\
v_{k+1} &= v_k + w_v \\
\omega_{k+1} &= \omega_k + w_\omega
\end{align}

In compact form

\begin{equation}
x_{k+1} = f(x_k) + w_k
\end{equation}

where

\begin{equation}
w_k =
\begin{bmatrix}
0 \\
0 \\
0 \\
w_v \\
w_\omega
\end{bmatrix}
\in \mathbb{R}^5
\end{equation}

\section*{Process Noise}

The process noise is modeled as Gaussian

\begin{equation}
w_k \sim \mathcal{N}(0,Q)
\end{equation}

with covariance

\begin{equation}
Q =
\begin{bmatrix}
0 &0&0&0&0\\
0 &0&0&0&0\\
0 &0&0&0&0\\
0 &0&0&\sigma_v^2&0\\
0 &0&0&0&\sigma_\omega^2
\end{bmatrix}
\end{equation}

\begin{center}
\begin{tabular}{lll}
\toprule
Parameter & Description & Units \\
\midrule
$\sigma_v^2$ & Variance of velocity noise & $(m/s)^2$ \\
$\sigma_\omega^2$ & Variance of yaw-rate noise & $(rad/s)^2$ \\
\bottomrule
\end{tabular}
\end{center}

\section*{Measurement Model}

Wheel encoders measure left and right wheel displacements

\begin{equation}
\Delta s_L, \quad \Delta s_R
\end{equation}

over timestep $\Delta t$.

The forward velocity and yaw rate are computed as

\begin{align}
v_{enc} &= \frac{\Delta s_R + \Delta s_L}{2 \Delta t} \\
\omega_{enc} &= \frac{\Delta s_R - \Delta s_L}{b \Delta t}
\end{align}

where

\begin{center}
$b$ = distance between the drive wheels (meters)
\end{center}

The measurement vector is

\begin{equation}
z_k =
\begin{bmatrix}
v_{enc} \\
\omega_{enc}
\end{bmatrix}
\in \mathbb{R}^2
\end{equation}

The measurement equation is linear

\begin{equation}
z_k = H x_k + n_k
\end{equation}

with

\begin{equation}
H =
\begin{bmatrix}
0 &0&0&1&0 \\
0 &0&0&0&1
\end{bmatrix}
\end{equation}

and measurement noise

\begin{equation}
n_k \sim \mathcal{N}(0,R), \quad R \in \mathbb{R}^{2 \times 2}
\end{equation}

\section*{Summary}

The resulting estimator uses

\begin{itemize}
\item State dimension: $5$
\item Measurement dimension: $2$
\item Nonlinear process model
\item Linear measurement model
\end{itemize}

This structure is well suited for an Extended Kalman Filter implementation
for differential drive robots using wheel encoder measurements.

\end{document}
